% ==========================================================================
% ch10.tex — dihasilkan otomatis dari website/chapters/ch10.qmd
% Boleh disunting manual. 'sync' hanya menimpa bila .qmd berubah DAN .tex
% belum disunting; kalau keduanya berubah, dilewati (pakai --force).
% qmd-sha1: 9357996c9227
% tex-sha1: 2e449f693a83
% ==========================================================================
% >>>>> ISI OTOMATIS (jangan hapus baris ini) >>>>>
\chapter{Deret Waktu \& Data Sensor}\label{deret-waktu-data-sensor}

\section{Konstruksi Sampel Temporal: Lookback, Stride, Horizon, dan Penyelarasan}\label{konstruksi-sampel-temporal-lookback-stride-horizon-dan-penyelarasan}

Dalam data tabular konvensional, satu baris umumnya merepresentasikan satu unit observasi yang independen. Namun, pada data deret waktu (\emph{time series}), setiap observasi terhubung oleh urutan kronologis. Baris-baris tersebut tidak dapat diambil secara acak. Untuk melatih algoritma pada data kontinu, sampel pembelajaran harus dikonstruksi melalui mekanisme \emph{windowing}.

Mekanisme \emph{windowing} mempartisi deret waktu kontinu menjadi himpunan sampel individual, di mana masing-masing memuat fitur historis sebagai masukan dan target masa depan sebagai prediksi. Proses ini dikendalikan oleh tiga parameter operasional: \emph{lookback}, \emph{horizon}, dan \emph{stride}. Parameter \emph{lookback} menentukan panjang rentang waktu historis yang digunakan model, misalnya tujuh hari riwayat penjualan untuk memprediksi hari esok. Parameter \emph{horizon} menentukan jarak waktu masa depan serta panjang target prediksi tersebut. Selanjutnya, \emph{stride} mengatur ukuran pergeseran jendela saat membentuk sampel berikutnya. Nilai \emph{stride} yang kecil akan memperbanyak jumlah sampel latih, tetapi sekaligus meningkatkan korelasi antarobservasi akibat tingginya tingkat tumpang tindih (\emph{overlapping}).

Secara matematis, untuk deret waktu \(x = \{x_1, x_2, \dots, x_T\}\), sampel historis pada titik waktu \(t\) direpresentasikan sebagai vektor jendela \(W_t = [x_{t-L+1}, x_{t-L+2}, \dots, x_t]\). Dalam formulasi ini, variabel \(W_t\) merupakan jendela input pada titik waktu \(t\) yang memiliki panjang \emph{lookback} sebesar \(L\). Jendela tersebut dipasangkan dengan target observasi \(y_{t+H}\), dengan variabel \(H\) menunjukkan jarak \emph{horizon} ke masa depan.

\begin{figure}[htbp]\centering
\includegraphics[width=\linewidth,height=0.72\textheight,keepaspectratio]{ch10-fig-1}
\caption{Jendela Rolling dan Expanding pada Deret Waktu}
\label{fig:ch10-fig-1}
\end{figure}

Penyelarasan waktu antara fitur historis dan target masa depan sangat rentan memicu \emph{data leakage}. Praktik \emph{pipeline} yang benar mengharuskan perancangan jendela waktu secara terstruktur. Pendekatan pertama, \emph{causal window}, menempatkan batas akhir jendela input tepat sebelum awal \emph{horizon}. Metode ini memastikan model hanya memproses data historis untuk memprediksi masa depan, menjadikannya syarat mutlak untuk tugas \emph{forecasting}. Pendekatan kedua, \emph{centered window}, merangkum observasi sebelum dan sesudah suatu titik waktu. Penggunaan skema ini pada kasus \emph{forecasting} akan mendistorsi model karena informasi masa depan masuk sebagai input. Oleh karenanya, \emph{centered window} dibatasi hanya untuk analisis retrospektif, seperti imputasi nilai kosong atau penghalusan tren.

Terakhir, pembentukan sampel ini menuntut interval waktu antarpengamatan yang konstan. Praktiknya, data sensor atau log transaksi sering tercatat dengan interval tidak beraturan sehingga wajib diseragamkan (\emph{resampling}) terlebih dahulu. Pengguna dapat menerapkan \emph{downsampling} untuk mengagregasi frekuensi tinggi menjadi interval yang lebih rendah, atau \emph{upsampling} untuk mengestimasi nilai pada celah antarpengamatan. Dengan interval yang konstan, setiap pergeseran \emph{stride} akan memuat rentang historis yang ekuivalen untuk dipelajari oleh algoritma.

\section{\texorpdfstring{Fitur Lag dan Statistik \emph{Rolling}}{Fitur Lag dan Statistik Rolling}}\label{fitur-lag-dan-statistik-rolling}

Setelah mengonstruksi \emph{lookback window}, langkah selanjutnya adalah mengekstraksi rentang waktu tersebut menjadi atribut diskrit. Tujuan utama tahap ini adalah meratakan dimensi waktu. Model \emph{machine learning} tabular standar beroperasi pada baris dan kolom yang saling lepas dan tidak memiliki mekanisme bawaan untuk mengenali urutan kejadian. Dengan memetakan riwayat pengamatan menjadi kolom-kolom baru, kita mengonversi proses historis yang dinamis menjadi deretan prediktor statis.

Cara paling langsung untuk merepresentasikan masa lalu adalah melalui fitur \emph{lag}. Fitur ini menarik nilai observasi dari langkah waktu sebelumnya untuk digunakan sebagai prediktor pada saat ini. Sebagai contoh, jika kita memprediksi suhu udara hari ini, suhu dari hari kemarin bertindak sebagai fitur \emph{lag} 1, sedangkan suhu dua hari yang lalu menjadi \emph{lag} 2.

Meskipun fitur \emph{lag} efektif menyediakan referensi terdekat, penggunaannya yang berlebihan dapat memunculkan kendala struktural. Ekstraksi riwayat yang terlalu jauh ke belakang secara terpisah akan menambah dimensi matriks \emph{dataset} secara eksponensial. Pelebaran matriks yang diikuti oleh korelasi antarfitur ini menyulitkan model dalam membedakan sinyal prediktif dari \emph{noise}, sehingga probabilitas \emph{overfitting} meningkat.

Sebagai alternatif untuk menangkap pola historis tanpa memperbanyak dimensi kolom, pengguna menerapkan agregasi statistik. Pendekatan pertama adalah menggunakan statistik \emph{rolling}, yang meringkas \emph{lookback window} berukuran konstan menjadi satu metrik agregat. Sebagai contoh, alih-alih mengekstraksi observasi 30 hari sebagai 30 prediktor terpisah, data tersebut dihitung menjadi nilai rata-rata, standar deviasi untuk mendeteksi deviasi fluktuasi, atau nilai puncak observasi. Metrik representatif dalam kelompok ini mencakup \emph{Simple Moving Average}, yang diformulasikan sebagai \(SMA_t = \frac{1}{w} \sum_{i=0}^{w-1} x_{t-i}\). Pada perumusan tersebut, \(SMA_t\) adalah nilai rata-rata bergerak pada waktu \(t\), variabel \(w\) menentukan ukuran jendela, dan \(x_{t-i}\) adalah nilai aktual pada \(i\) langkah sebelum \(t\). Penentuan ukuran jendela ini membutuhkan keseimbangan fungsional; jendela yang sempit merekam dinamika terkini namun sangat terpengaruh oleh \emph{noise}, sedangkan jendela yang panjang menyediakan stabilitas pembacaan tren namun lambat dalam merefleksikan perubahan arah musiman.

Pendekatan kedua adalah menggunakan statistik \emph{expanding}. Berbeda dengan pergeseran interval pada metode \emph{rolling}, \emph{expanding window} menahan titik awal pengamatan pada batas tetap, sedangkan ujung observasi berlanjut sejalan masuknya baris \emph{dataset} baru. Mekanisme ini diaplikasikan untuk membangun rekam jejak kumulatif masa lalu, seperti metrik total volume transaksi sistem sejak fase aktivasi awal hingga waktu observasi terkini.

\begin{figure}[htbp]\centering
\includegraphics[width=\linewidth,height=0.72\textheight,keepaspectratio]{ch10-fig-2}
\caption{Perbandingan Mekanisme Jendela Bergeser (Rolling) vs Melebar (Expanding)}
\label{fig:ch10-fig-2}
\end{figure}

Melalui kombinasi tiga teknik representasi tersebut (referensi jarak dekat dari \emph{lag}, deteksi tren lokal dari \emph{rolling}, serta observasi kumulatif dari \emph{expanding}), karakteristik data temporal dipetakan secara terstruktur ke dalam format tabular. Pada lingkungan operasional, pustaka pemrosesan runtun waktu seperti \emph{tsfresh} atau komponen transformasi \emph{aeon} melakukan perhitungan puluhan agregasi statistik ini secara sistematis.

\section{Differencing, Tren, dan Pola Musiman}\label{differencing-tren-dan-pola-musiman}

Data deret waktu memiliki dua pola makro yang membedakannya dari data tabular biasa: tren dan musiman. Tren adalah pergerakan arah jangka panjang, seperti angka penjualan yang terus naik dari tahun ke tahun. Musiman adalah fluktuasi periodik yang berulang secara teratur, misalnya lonjakan pengunjung situs setiap akhir pekan.

Bentuk makro seperti tren memberikan hambatan komputasi bagi algoritma, khususnya pada kelompok model berbasis pohon keputusan seperti \emph{random forest} atau \emph{gradient boosting}. Arsitektur klasifikasi ini mempartisi distribusi fitur berdasarkan rentang observasi pelatihan tanpa kemampuan ekstrapolasi. Jika data operasional di lapangan mencapai nilai yang melampaui batas maksimum dari data pelatihan, model tersebut tidak mampu memprediksi titik puncak baru tersebut meskipun observasi secara proporsional memperlihatkan lintasan tren yang meningkat.

Untuk menengahi keterbatasan ekstrapolasi ini, representasi data perlu dikonversi mencapai ekuilibrium stasioner. Stasioneritas merupakan kondisi kestabilan statistik di mana metrik seperti rata-rata dan varians tidak lagi mengalami pergeseran seiring berjalannya waktu. Mekanisme utama untuk membentuk status stasioner adalah transformasi \emph{differencing}. Pendekatan ini mengubah fokus algoritma dari memprediksi nilai akumulatif menjadi memprediksi jarak diferensiasi antar titik waktu konsekutif.

\begin{figure}[htbp]\centering
\includegraphics[width=\linewidth,height=0.72\textheight,keepaspectratio]{ch10-fig-3}
\caption{Ekstraksi Fitur Menggunakan tsfresh yang Meratakan Jendela Deret Waktu ke Tabular}
\label{fig:ch10-fig-3}
\end{figure}

Secara matematis, proses \emph{differencing} orde pertama dikalkulasikan melalui rumusan \(y'_t = y_t - y_{t-1}\). Dalam persamaan tersebut, \(y'_t\) merujuk pada nilai luaran hasil transformasi pada langkah waktu \(t\). Adapun \(y_t\) merepresentasikan matriks observasi sebelum modifikasi, dan \(y_{t-1}\) merupakan variabel kondisi observasi dari posisi satu waktu sebelumnya. 

Praktik rekayasa fitur menyediakan ragam modifikasi lain dalam menangani tren maupun periodisitas. Proses \emph{seasonal differencing} digunakan untuk mereduksi efek repetisi dengan menghitung diferensiasi nilai posisi terbaru dengan lokasi ekuivalen pada rotasi sebelumnya, yaitu \(y'_t = y_t - y_{t-m}\), tempat parameter \(m\) membatasi durasi jarak antarpola (misalnya panjang \(m=7\) untuk aktivitas mingguan). Teknik dekomposisi deret waktu juga dapat dieksekusi melalui prosedur analitis STL (\emph{Seasonal and Trend decomposition using Loess}) untuk menyaring sinyal observasi menjadi fraksi tren, sisaan residu, dan efek musim; lalu melatih prediksi model terfokus hanya pada sisaan komponen stasioner residu. Alternatif transformasi lainnya berfokus untuk mencatat pergerakan melalui ekstraksi kemiringan sudut (\emph{slope regression}) dalam blok waktu, maupun penambahan kolom fitur elemen kalender (indikator kategorikal waktu). Penambahan label musim ini mendampingi analisis pola sehingga parameter tren orisinal tetap berada pada format asalnya.

\section{Fitur Domain Frekuensi}\label{fitur-domain-frekuensi}

Seluruh teknik ekstraksi yang memproses urutan nilai berdasarkan waktu, seperti \emph{lag}, \emph{differencing}, maupun statistik \emph{rolling}, beroperasi di dalam domain waktu. Pendekatan ini mengekstrak pola observasi melalui analisis perubahan amplitudo secara berurutan. Format \emph{feature} berbasis waktu tersebut sangat efektif dalam memodelkan set data berfrekuensi rendah, contohnya perubahan agregat penjualan harian atau indikator fluktuasi ekonomi bulanan, di mana transisi absolut pada titik parameter terlihat konstan.

Namun, keterbatasan parameter domain waktu mulai terlihat saat mengelola sinyal data sensor berfrekuensi tinggi. Rangkaian variabel fisik berosilasi rapat, seperti transmisi akselerometer di gawai, sensor getaran di peralatan manufaktur, ataupun pembacaan elektrokardiogram (ECG), mencatat ribuan pengukuran dalam waktu yang sangat singkat. Pada tingkat frekuensi ini, observasi pergeseran simpangan secara berurutan per milidetik jarang merepresentasikan sinyal prediktif yang solid. Pola karakteristik pembentuknya lebih dominan berada pada rasio perulangan gelombang. Untuk mengakomodasi dinamika struktur periodik ini, fokus ekstraksi harus digeser menuju domain frekuensi.

Transisi analitis tersebut dieksekusi secara primer menggunakan metode matematis transformasi Fourier, khususnya algoritma operasional \emph{Fast Fourier Transform} (FFT). Secara fungsi konseptual, perhitungan FFT memisahkan jendela sinyal gabungan ke dalam struktur turunan dari beberapa gelombang komponen sederhana, di mana setiap pecahan gelombang dikategorikan berdasarkan rasio perulangan dan amplitudonya.

\begin{figure}[htbp]\centering
\includegraphics[width=\linewidth,height=0.72\textheight,keepaspectratio]{ch10-fig-4}
\caption{Skema Purge Gap dalam Pemisahan Data Temporal}
\label{fig:ch10-fig-4}
\end{figure}

Penerapan transformasi konseptual ini mendefinisi ulang arsitektur \emph{feature}. Menggantikan pencatatan amplitudo sinyal pada sumbu linier, observasi dikonversi menjadi perhitungan spektrum frekuensi utama (contohnya keberadaan dominan gelombang 50 Hz) dalam keseluruhan batas parameter jendela. Transposisi kerangka waktu menjadi distribusi frekuensi diformulasikan secara teknis menggunakan \emph{Discrete Fourier Transform} (DFT), yang memiliki persamaan dasar \(X_k = \sum_{n=0}^{N-1} x_n \cdot e^{-i 2 \pi k n / N}\). Dalam persamaan tersebut, \(X_k\) adalah komponen spektrum kompleks untuk frekuensi ke-\(k\) yang memuat informasi amplitudo dan fase. Adapun \(x_n\) adalah nilai observasi pada titik waktu ke-\(n\), dan \(N\) adalah jumlah total observasi dalam satu blok perhitungan. Faktor eksponensial \(e^{-i 2 \pi k n / N}\) berperan sebagai gelombang acuan yang mengukur seberapa kuat frekuensi \(k\) hadir dalam sinyal.

Dalam praktiknya, keluaran mentah FFT jarang langsung dipakai untuk pemodelan tabular. Yang diekstrak justru beberapa ringkasan dari spektrum. Puncak frekuensi (\emph{peak frequency}) menandai frekuensi dengan amplitudo terkuat. \emph{Band energy} mengukur total daya pada rentang frekuensi tertentu, sedangkan \emph{spectral centroid} menunjukkan pusat massa sebaran daya di sepanjang spektrum. Entropi spektral melengkapi ukuran ini: nilai yang tinggi menandakan sinyal menyebar seperti \emph{noise}, sedangkan nilai yang rendah menandakan daya terpusat pada sedikit frekuensi.

Dalam \emph{pipeline} rekayasa fitur modern, pengguna tidak perlu membangun algoritma komputasi spektral secara mandiri karena ekosistem pustaka \emph{machine learning} telah menyediakan kerangka kerja fungsional. Modul pemrosesan seperti \passthrough{\lstinline!tsfresh!} dirancang untuk menghitung puluhan variabel frekuensi, mencakup metrik koefisien FFT, kerapatan spektral Welch, hingga entropi komputasi Fourier pada rentang jendela \emph{rolling}. Melengkapi instrumen otomatisasi tersebut, sistem spesifik \emph{pipeline} seperti \passthrough{\lstinline!aeon!} mengintegrasikan fungsionalitas \emph{transformer} data seperti \passthrough{\lstinline!PeriodogramTransformer!} maupun \passthrough{\lstinline!DWTTransformer!} yang memodifikasi komponen operasional \emph{wavelet} ke dalam rancangan prediksi yang bebas dari \emph{data leakage}.

Salah satu aplikasi terstruktur dari representasi frekuensi adalah pada infrastruktur pemeliharaan prediktif (\emph{predictive maintenance}) sistem manufaktur mekanik. Kerusakan awal pada komponen bantalan poros mesin (\emph{bearing}) jarang terdeteksi lewat volume getaran keseluruhan; kondisi aus tersebut justru terekam pada pergeseran parameter fluktuatif serta kenaikan daya spektrum di level tinggi secara sangat terlokalisasi. Ekstraksi spesifik \emph{feature} dari domain spektral inilah yang membuat fungsi pemantauan sanggup mendeteksi deviasi anomali sejak fase awal yang menjadi peringatan dini pencegahan sebelum perangkat memburuk secara fisik. Pembelajaran mengenai mekanisme dekomposisi tersebut turut membangun kerangka teknis analitis gelombang berbasis akustik di Bab 12 yang mengubah aliran suara audio ke dalam matriks gambar dua dimensi berformat spektrogram.

\section{Validasi Temporal: Menghindari Look-Ahead Leakage}\label{validasi-temporal-menghindari-look-ahead-leakage}

Penerapan \emph{cross-validation} acak seperti \emph{K-Fold} pada data deret waktu akan memicu \emph{data leakage}. Pengacakan baris sebelum membagi \emph{dataset} tidak memedulikan susunan kronologis. Sebagai ilustrasi, jika model dilatih menggunakan pengamatan hari Jumat untuk memprediksi target pada hari Kamis, informasi masa depan tersebut sudah diserap oleh fase pelatihan dan validasi. Skema ini menghasilkan evaluasi akurasi semu, di mana performa aktual akan menurun drastis saat algoritma memproses aliran data baru. Pemisahan data untuk kasus temporal wajib mematuhi pemotongan hierarki kronologis (\emph{temporal split}), menempatkan alokasi pelatihan secara eksklusif sebelum rentang blok pengujian. Desain ini menggunakan dua varian; pertama, \emph{expanding window} yang memperbesar ukuran jendela pelatihan namun mempertahankan rentang ukuran uji; dan kedua, \emph{sliding window} di mana blok parameter set data bergeser konstan mengkuti pergerakan kalender operasional.

Meskipun mekanisme \emph{temporal split} menjaga isolasi urutan parameter waktu, formasi observasi berpotensi mengintroduksi bias melalui persinggungan data. Operasional \emph{windowing} menstruktur \emph{dataset} temporal ke dalam barisan parameter yang membentuk \emph{overlapping}. Kondisi ini menyebabkan observasi terakhir di set pelatihan dan \emph{dataset} pertama pada modul validasi bisa saja menggunakan sumber elemen observasi masa lalu yang ekuivalen pada parameter \emph{lookback}. Misalnya pada pemodelan dengan batas \emph{lookback} tujuh hari, baris validasi perdana pada rentang hari ke-100 merekonstruksi rentang matriks mentah dari nilai ke-93 hingga 99; sedangkan komponen data mentah serupa masih masuk ke dalam bagian ekstraksi pengujian observasi hari ke-99 yang dialokasikan di dalam ranah pelatihan. Garis demarkasi perpotongan antar-\emph{dataset} yang bersinggungan langsung tersebut mengaburkan batas objektivitas evaluasi prediksi.

\begin{figure}[htbp]\centering
\includegraphics[width=\linewidth,height=0.72\textheight,keepaspectratio]{ch10-fig-5}
\caption{Letak Purge Gap di Antara Set Pelatihan Terakhir dan Set Validasi Pertama}
\label{fig:ch10-fig-5}
\end{figure}

Praktik \emph{pipeline} yang benar mengatasi kerentanan \emph{data leakage} di garis batas pemotongan data dengan cara menyisipkan \emph{purge gap}. Pemisahan \emph{purge gap} bertindak sebagai alokasi observasi perantara yang dieliminasi secara terencana untuk menghindari perpindahan data residu antara parameter batas observasi. Penghitungan panjang observasi yang dibuang diukur melalui parameter eksak \(n_{\text{purge}} = w_{\text{lookback}} + h_{\text{horizon}}\). Formula matematika ini menggunakan variabel \(n_{\text{purge}}\) untuk mewakili ukuran celah data yang dieksklusi, serta \(w_{\text{lookback}}\) sebagai indeks blok langkah historis ke arah belakang. Tambahan variabel \(h_{\text{horizon}}\) mengindikasikan ukuran \emph{horizon} atau unit masa depan prediksi observasi. Parameter prediksi masa depan wajib masuk ke dalam operasi penambahan tersebut karena label target pada rentetan sampel pelatihan ujung merujuk ke elemen data terdepan. Dengan meloncati segmen ini, fungsi pengujian di dalam observasi validasi memperoleh status keterpisahan mutlak tanpa anomali duplikasi elemen data. Dalam sistem \emph{machine learning} saat ini, kerangka fungsional pemrosesan prediksi seperti \passthrough{\lstinline!sktime!} mengoperasikan parameter internal untuk secara spesifik membuang area persinggungan data dari prosedur \emph{cross-validation}. Inisialisasi ruang kosong ini menjamin data \emph{holdout} di fase validasi memperoleh ketahanan evaluasi prediksi.

\section{Fitur untuk Model Klasik vs.~Input Sekuensial}\label{fitur-untuk-model-klasik-vs.-input-sekuensial}

Bentuk rekayasa fitur temporal bergantung pada jenis model yang akan mengolahnya. Model berbasis pohon seperti \emph{random forest} atau \emph{gradient boosting} menerima input sebagai vektor satu dimensi, sehingga urutan waktu harus diratakan lebih dahulu menjadi kolom-kolom terpisah. Model semacam ini tidak menyimpan informasi urutan: setiap fitur dievaluasi sendiri-sendiri tanpa mengetahui posisinya dalam waktu. Karena itu, dimensi waktu perlu dikodekan secara eksplisit melalui fitur turunan seperti \emph{lag} dan agregasi statistik. Pustaka seperti \passthrough{\lstinline!tsfresh!} mengotomasi langkah ini dengan menghitung ribuan fitur sekaligus, lalu meringkasnya menjadi tabel dua dimensi yang siap dilatih.

\begin{figure}[htbp]\centering
\includegraphics[width=\linewidth,height=0.72\textheight,keepaspectratio]{ch10-fig-6}
\caption{Skema Validasi Temporal dengan Pemisahan Kronologis dan Area Purge Gap}
\label{fig:ch10-fig-6}
\end{figure}

Model sekuensial seperti RNN, LSTM, dan \emph{transformer} bekerja sebaliknya. Fitur tidak diratakan, tetapi dipertahankan sebagai \emph{tensor} dua dimensi dengan satu sumbu langkah waktu dan satu sumbu fitur. Struktur ini menjaga urutan kronologis: model menghubungkan observasi masa lalu dengan masa depan lewat mekanisme internalnya, yaitu \emph{hidden state} pada RNN atau \emph{attention} pada \emph{transformer}. Dengan begitu, ekstraksi pola temporal tidak lagi dirancang manual, melainkan dipelajari sendiri sebagai bobot model selama \emph{training}.

Meski jaringan saraf mampu mengenali urutan observasi, data mentahnya tetap perlu disiapkan, terutama diskalakan agar \emph{gradient} stabil selama \emph{training}. Model ini juga masih memerlukan \emph{cyclical encoding} untuk merepresentasikan waktu yang berulang. \emph{Encoding} tersebut memetakan waktu ke lingkaran melalui fungsi sinus dan kosinus, \(x_{sin} = \sin\left(\frac{2\pi t}{P}\right)\) dan \(x_{cos} = \cos\left(\frac{2\pi t}{P}\right)\), dengan \(t\) sebagai posisi waktu dan \(P\) sebagai panjang satu siklus penuh, misalnya 24 untuk jam dalam sehari. Tanpa \emph{encoding} ini, model menganggap jam 23 dan jam 0 berjauhan, padahal keduanya berdampingan.

Pendekatan modern mulai menggabungkan kedua gaya ini. Lag-Llama, misalnya, memadukan fitur \emph{lag} dengan input mentah dalam satu model. PatchTST memecah deret waktu menjadi potongan-potongan (\emph{patch}) pendek sebelum diproses, sehingga korelasi antarsegmen lebih mudah ditangkap. Ada pula model \emph{foundation} khusus deret waktu seperti \emph{Chronos} dan MOMENT yang dipakai sebagai \emph{feature extractor}: model beku ini mengubah deret mentah menjadi vektor \emph{embedding} padat, yang kemudian diumpankan ke model klasik pada tahap berikutnya.

\section{Studi Kasus: Prediksi Deret Waktu Multivariat}\label{studi-kasus-prediksi-deret-waktu-multivariat}

Studi kasus ini menerapkan rekayasa fitur deret waktu multivariat untuk memprediksi beban listrik suatu wilayah pada langkah waktu berikutnya. Beban listrik dipengaruhi banyak faktor sekaligus, sehingga suhu harian, kelembapan lokal, dan volume konsumsi dicatat sebagai fitur masukan model multivariat. Sebagai \emph{baseline}, kita gunakan \emph{random forest}. Karena model ini menerima input tabular, deret waktu harus diratakan lebih dahulu menjadi baris-baris tunggal. Kita tetapkan \emph{horizon} sepanjang satu langkah ke depan sebagai target, dan \emph{lookback} sepanjang 24 langkah agar setiap sampel merangkum pola beban satu siklus harian sebelumnya.

Dari jendela waktu itu, kita bentuk beberapa kelompok fitur. Pertama, jam (0--23) dikodekan sebagai \emph{cyclical features} dengan sinus dan kosinus, agar model tidak menganggap jam 0 dan jam 23 berjauhan padahal berdampingan. Kedua, 24 fitur \emph{lag} beban listrik menyalin nilai satu hingga 24 jam sebelumnya untuk menangkap pola harian. Ketiga, fitur \emph{rolling} meringkas tren jangka pendek, misalnya rata-rata bergerak suhu, \(MA_t = \frac{1}{w} \sum_{i=0}^{w-1} x_{t-i}\), dengan ukuran jendela \(w=24\) yang merata-ratakan suhu 24 jam terakhir dari titik prediksi.

Tahap selanjutnya adalah perancangan validasi evaluasi prediksi. Proses pengocokan \emph{dataset} baris matriks sangat tidak disarankan karena modifikasi ini mendistorsi kerangka hierarki sekuensial yang menyebabkan masalah \emph{data leakage}. Sebagai pengganti, prosedur evaluasi dijalankan berbekal \emph{temporal split} eksklusif, contohnya memisahkan blok observasi kronologis berupa penggunaan matriks pengamatan murni sepuluh bulan awal dan merelokasi sisa parameter \emph{dataset} berikutnya ke kelompok parameter fungsi pengujian akhir.

Meski pemisahan sudah kronologis, sampel uji pertama tetap membawa 24 fitur \emph{lag} yang nilainya diambil dari sisi pelatihan, tepat di seberang garis potong. Akibatnya, sampel pelatihan terakhir dan sampel uji pertama berbagi sebagian data mentah yang sama. Untuk menghapus irisan ini, kita sisipkan \emph{purge gap} di antara keduanya.

\begin{lstlisting}[language=Python]
import pandas as pd
import numpy as np

# 1. Menyiapkan data deret waktu sintetis (frekuensi per jam selama 30 hari)
date_range = pd.date_range(start='2026-01-01', periods=24*30, freq='h')
np.random.seed(42)
df = pd.DataFrame({
    'Suhu': np.random.normal(loc=25, scale=5, size=len(date_range)),
    'Beban_Listrik': np.random.normal(loc=150, scale=30, size=len(date_range))
}, index=date_range)

# 2. Menentukan titik potong temporal (misalnya: 80% pelatihan, 20% pengujian)
split_point = df.index[int(len(df) * 0.8)]

# 3. Mendefinisikan ukuran Purge Gap (contoh: 25 jam)
purge_gap_hours = 25
purge_gap_delta = pd.Timedelta(hours=purge_gap_hours)

# 4. Melakukan pemisahan temporal secara ketat dengan menyertakan Purge Gap
X_train = df[:split_point]
# Set pengujian dimulai setelah split_point + purge_gap_delta untuk membuang overlap
X_test = df[split_point + purge_gap_delta:]

# 5. Memverifikasi isolasi temporal
print(f"Batas Terakhir Set Pelatihan: {X_train.index[-1]}")
print(f"Batas Awal Set Pengujian   : {X_test.index[0]}")
print(f"Jarak Waktu Antar Set      : {X_test.index[0] - X_train.index[-1]}")
print(f"Jumlah Baris Terbuang (Gap): {len(df) - (len(X_train) + len(X_test))} jam")
\end{lstlisting}

Kode di atas menjalankan bagian inti dari isolasi temporal, yaitu penyisipan \emph{purge gap}. Celah sepanjang 25 jam dipotong tepat di titik pemisahan; angka ini setara dengan panjang \emph{lookback} ditambah \emph{horizon}, sehingga cukup untuk membuang seluruh irisan data antara set pelatihan dan pengujian. Pada \emph{pipeline} yang lengkap, langkah ini dilanjutkan dengan penskalaan yang di-\emph{fit} hanya pada data pelatihan lalu diterapkan ke data uji, agar statistik masa depan tidak ikut membentuk parameter praproses. Dengan jendela yang dirancang rapi dan \emph{purge gap} yang memisahkan kedua set, metrik yang dihasilkan mencerminkan kemampuan model di data baru sebelum diterapkan pada skala \emph{production}.
